\PassOptionsToPackage{dvipsnames,svgnames,table}{xcolor}
\documentclass[11pt]{article}
\usepackage[a4paper,margin=22mm]{geometry}
\usepackage[no-math]{fontspec}
\setmainfont{DejaVu Serif}
\usepackage{amsmath,amssymb,amsthm,mathtools,graphicx,xcolor,hyperref,xurl,xspace,ifthen,enumerate}
\usepackage[shortlabels]{enumitem}
\setlength{\emergencystretch}{4em}
\AtBeginDocument{\special{pdf:mapline pzdr Dingbats <uzdr.pfb}}
\SetMathAlphabet{\mathbf}{normal}{OT1}{cmr}{bx}{n}
\SetMathAlphabet{\mathbf}{bold}{OT1}{cmr}{bx}{n}
\newtheorem{theo}{Theorem}[section]
\newtheorem{lemm}[theo]{Lemma}
\newtheorem{thm}[theo]{Theorem}
\newtheorem{proposition}[theo]{Proposition}
\newtheorem{corollary}[theo]{Corollary}
\newtheorem{definition}[theo]{Definition}
\newtheorem{theorem}[theo]{Theorem}
\newtheorem*{remas*}{Remarks}
\newtheorem{remas}[theo]{Remarks}
\newtheorem*{exam*}{Example}
\newtheorem{ques}[theo]{Question}
\providecommand{\og}{«\nobreakspace}
\providecommand{\fg}{\nobreakspace»}
\newtheorem{lem}[theo]{Lemma}
\newtheorem{lemma}[theo]{Lemma}
\newtheorem{prop}[theo]{Proposition}
\newtheorem{coro}[theo]{Corollary}
\newtheorem{corol}[theo]{Corollary}
\newtheorem*{conj*}{Conjecture}
\newtheorem{conj}[theo]{Conjecture}
\newtheorem{conjecture}[theo]{Conjecture}
\newtheorem{Proposition}[theo]{Proposition}
\newtheorem{theoreme}[theo]{Theorem}
\theoremstyle{definition}
\newtheorem{defi}[theo]{Definition}
\newtheorem{exam}[theo]{Example}
\newtheorem{exem}[theo]{Example}
\newtheorem{remark}[theo]{Remark}
\newtheorem{example}[theo]{Example}
\newtheorem{exams}[theo]{Examples}
\newtheorem{rema}[theo]{Remark}
\newtheorem{nota}[theo]{Notation}
\newtheorem{notation}[theo]{Notation}
\newtheorem*{notas}{Notations}
\newtheorem{result}[theo]{Result}
\newtheorem{question}[theo]{Question}
\newtheorem{prob}[theo]{Problem}
\newtheorem*{rema*}{Remark}
\newtheorem*{defi*}{Definition}
\newtheorem*{theo*}{Theorem}
\newtheorem*{lemm*}{Lemma}
\newtheorem*{prop*}{Proposition}
\newcommand{\enoncename}{}
\newtheorem{corpusenonce}[theo]{\enoncename}
\newenvironment{enonce}[2][]{\renewcommand{\enoncename}{#2}\begin{corpusenonce}}{\end{corpusenonce}}
\newenvironment{enonce*}[2][]{\par\medskip\noindent\textbf{#2.}\itshape}{\par\medskip}
\newenvironment{proofc}[1][\proofname]{\begin{proof}[#1]}{\end{proof}}
\newcommand{\Mc}{\mathcal{M}}
\newcommand{\pointir}{.\ ---\ }
\newcommand{\equalenv}[2]{\expandafter\let\csname#1\expandafter\endcsname\csname#2\endcsname\expandafter\let\csname end#1\expandafter\endcsname\csname end#2\endcsname}
\providecommand{\diagbox}[3][]{\begin{tikzpicture}[x=1em,y=1em,baseline=(current bounding box.center)]\useasboundingbox(0,0) rectangle(3,1.4);\draw(0,1.4)--(3,0);\node at(.5,.3){#2};\node at(2.5,1.1){#3};\end{tikzpicture}}
\providecommand{\eg}{e.g.\ }
\providecommand{\ie}{i.e.\ }
\providecommand{\cf}{cf.\ }
\providecommand{\longthanks}[1]{\par\smallskip\noindent #1\par\smallskip}
\providecommand{\qbinom}[2]{\genfrac{[}{]}{0pt}{}{#1}{#2}}

\providecommand{\mathof}[1]{\mathrm{#1}}

\numberwithin{equation}{section}


\def\endappendix{}


\usepackage{latexsym,amsmath,amssymb,verbatim,tikz,tikz-cd}
\usepackage{mathtools}
\usepackage{graphicx,xcolor}

\usepackage[english]{babel}
\usepackage{amssymb}
\usepackage{amsmath}
\usepackage{amsthm}
\usepackage{graphicx}
\usepackage{csquotes}
\usepackage{bm}
\usepackage{enumitem}
\usepackage{subcaption}
\usepackage{caption}
\renewcommand{\thesubfigure}{\Alph{subfigure}}
\usepackage{dirtytalk}
\providecommand{\nopunct}{\spacefactor1007\relax}
\newtheorem{problem}[theo]{Problem}

\newcommand{\sym}{\mathrm{Sym}}
\newcommand{\qsym}{\mathrm{QSym}}
\newcommand{\Des}{\operatorname{Des}}
\newcommand{\Asc}{\operatorname{Asc}}
\newcommand{\Syt}{\operatorname{SYT}}
\newcommand{\Short}{\operatorname{Short}}
\newcommand{\Loop}{\operatorname{Loop}}
\newcommand{\Stable}{\operatorname{Stable}}
\newcommand{\len}{\operatorname{len}}
\newcommand{\height}{\operatorname{height}}
\newcommand{\core}{\operatorname{core}}
\newcommand{\res}{\operatorname{res}}
\newcommand{\row}{\operatorname{row}}
\newcommand{\insrt}{\operatorname{insert}}
\newcommand{\UP}{\operatorname{UP}}
\newcommand{\DOWN}{\operatorname{DOWN}}


\newcommand{\ZZ}{\mathbb{Z}}
\newcommand{\NN}{\mathbb{N}}
\newcommand{\RR}{\mathbb{R}}
\newcommand{\QQ}{\mathbb{Q}}
\newcommand{\PP}{\mathbb{P}}

\newcommand{\F}{\mathcal{F}}
\newcommand{\Q}{\mathcal{Q}}
\newcommand{\A}{\mathcal{A}}
\newcommand{\M}{\mathcal{M}}
\newcommand{\I}{\mathcal{I}}
\newcommand{\PM}{\mathcal{PM}}
\newcommand{\Pcal}{\mathcal{P}}


\newcommand{\odds}{\mathrm{odds}}

\makeatletter
\newcommand{\proofstep}[1]{
  \par
  \addvspace{\medskipamount}
  \textbf{#1\@addpunct{.}}\enspace\ignorespaces
}

\newcommand\restr[2]{{
  \left.\kern-\nulldelimiterspace 
  #1 
  \vphantom{\big|} 
  \right|_{#2} 
  }}
\makeatother












\usepackage{cleveref}
\crefname{originaltheorem}{original theorem}{original theorems}
\Crefname{originaltheorem}{Original theorem}{Original theorems}
\makeatletter
\newlabel{thm:criterion schur positive young}{{2.5}{source}{}{sourcepointer-1}{}}
\newlabel{thm:criterion schur positive young@cref}{{[originaltheorem][0][]2.5}{[1][0][]source}{}{}{}}
\makeatother
\title{A set symmetric for a sparse statistic is Schur-positive with an explicit two-row expansion}
\author{Avichai Marmor}
\date{doi:10.5802/alco.351; CC BY 4.0}
\begin{document}
\maketitle
\noindent\textit{Editorial note.} The counted statement is the author's labelled lemma `lem:if two rows and symmetric then schur positive` (source lines 612-618); the proof environment begins on the next line after \textbackslash{}end\{lemma\}, so the pairing is by adjacency.\par\medskip
\noindent\textit{Editorial note.} Three context spans are retained: the author's Schur-positivity criterion, and the two superstandard characterisations that the proof invokes; every label used inside the retained text is therefore defined inside the package.\par\medskip
\noindent\textit{Editorial note.} The retained preamble is the editorially prepared standalone preamble retained for this source; packages used only elsewhere in the article are left out because no retained line needs them.\par\medskip
\noindent\textit{Editorial note.} No retained line refers to a label, citation or figure outside the retained spans, so the package declares no external reference.\par\medskip
\noindent\textit{Editorial scope.} One original result of the pinned publisher article, counted once: Schur-positivity of a set symmetric for a sparse statistic with its two-row expansion (author label lem:if two rows and symmetric then schur positive).\par\medskip
\noindent\textit{Allowed standard prerequisites (editorial).}\begin{itemize}
\item Symmetric set with respect to a statistic. A is symmetric for D : A -> 2\textasciicircum{}\{[N-1]\} if the generating function sum over a in A of t\textasciicircum{}\{D(a)\} expands in standard Young tableaux with coefficients depending only on the shape.. It is the hypothesis that produces equation (1) at source lines 620-623.
\item Sparse statistic. D(a) is sparse if it contains no two consecutive integers.. It is used at source line 634 to rule out shapes with more than two parts.
\item Superstandard tableau characterisations. The descent set of T\_lambda determines the dominance order of the shape, and in the two-row case it determines the tableau and the shape uniquely.. Lemma 1 at source lines 632-634 and Lemma 2 at source line 641 are the two inputs that turn the generating-function identity into the coefficient formula.
\end{itemize}\medskip
\ifdefined\hypertarget\hypertarget{sourcepointer-1}{}\fi
\noindent\textit{Source reference.} Original theorem 2.5: sources/alco-351-matching-schur-2024/publisher-source.tex, source lines 368--368. Original criterion of Adin and Roichman cited by the retained proof at source lines 620 and 632. It is a navigation pointer: the counted lemma's own proof is retained line for line and nothing is replaced by this reference. The reference number was inferred by counting the theorem-class environments of that section of the pinned source; it is not confirmed by the retained bibliography or by a label-number dump of the original, so it must be read as a navigation pointer only.\par\medskip
The power of these notions may be reflected by the following statement:
\begin{lemma}
\label{lem:superstandard characterize coefficients}
    Let $\lambda, \mu \vdash N$ be two partitions. Then:
    \begin{enumerate}
        \item $\Des(T_\lambda) = [N-1] \setminus \{\lambda_{1}', \lambda_{1}'+\lambda_{2}', \dots\}$.
        \item If $\Des(T) = [N-1] \setminus \{\lambda_{1}', \lambda_{1}'+\lambda_{2}', \dots\}$ for some $T \in \Syt(\mu)$, then $\mu' \ge \lambda'$. Furthermore, if $\mu = \lambda$ then $T = T_\lambda$.
    \end{enumerate}
\end{lemma}


Lemma~\ref{lem:superstandard characterize coefficients} associates the Young diagram of shape $\lambda$ with the set $\Des(T_\lambda)$. As we will see later, this association is powerful in analysing the Schur coefficients of symmetric sets.


Standard Young tableaux with at most $2$ rows have a slightly stronger property:
\begin{lemma}
\label{lem:superstandard two rows characterize coefficients}
    Let $\lambda = (N-k_1,k_1)$ and $\mu = (N-k_2,k_2)$ be two partitions of $N$ with at most two parts each. Then:
    \begin{enumerate}
        \item $\Des(T_\lambda) = \odds(k_1)$.
        \item If $\Des(T) = \odds(k_1)$ for some $T \in \Syt(\mu)$, then $k_1 = k_2$ and $T = T_\lambda$.
    \end{enumerate}
\end{lemma}

\begin{lemma}
\label{lem:if two rows and symmetric then schur positive}
    Let $\A$ be a symmetric set with respect to a sparse statistic $D:\A \to 2^{[N-1]}$, and denote $n = \lfloor\frac{N}{2} \rfloor$. Then $\A$ is Schur-positive, and its Schur expansion is
    \[
        \Q_{D} (\A) = \sum_{k=0}^n |\A(D = \odds(k))|\; s_{N-k,k}.
    \]
\end{lemma}

\begin{proof}
    The set $\A$ is assumed to be symmetric, so by Theorem~\ref{thm:criterion schur positive young}, we have
    \begin{equation}\label{eqn:proof symmetric set with no consecutive indices is schur positive}
        \sum_{a \in \A} \boldsymbol{t}^{D(a)} = \sum_{\lambda \vdash N} c_\lambda \sum_{T \in \Syt(\lambda)} \boldsymbol{t}^{\Des(T)},
    \end{equation}
    where $c_\lambda$ are the Schur coefficients of $\A$. It suffices to show that if there exists $k$ such that $\lambda=(N-k,k)$ then $c_\lambda = |\A(D = \odds(k))|$, and otherwise $c_\lambda=0$.

    If $\A = \emptyset$, then $\Q_{D}(\A) = 0$, and the statement holds. Therefore, we may assume that $\A \ne \emptyset$, and consequently, there exists $c_\lambda \ne 0$ for some partition $\lambda \vdash N$. Let $\mu \vdash N$ be a partition with $c_\mu \ne 0$, maximal in the conjugate order (\ie such that $c_\lambda = 0$ for every $\lambda \vdash N$ with $\lambda' > \mu'$).
    
    Equation~\eqref{eqn:proof symmetric set with no consecutive indices is schur positive} implies that the equation
    \begin{equation}\label{eqn:proof symmetric set with no consecutive indices is schur positive - specific set}
        |\A(D = J)| = \sum_{\lambda \vdash N} c_\lambda \left|\Syt(\lambda)(\Des = J)\right|
    \end{equation}
    holds for all $J \subseteq [N-1]$. Let us find the right-hand side of Equation~\eqref{eqn:proof symmetric set with no consecutive indices is schur positive - specific set} when substituting $J = \Des(T_\mu)$. Let $T \in \Syt(\lambda)$ be a tableau with $\Des(T) = \Des(T_\mu)$. By Lemma~\ref{lem:superstandard characterize coefficients}, we have $T = T_\mu$ or $\lambda' > \mu'$. However, if $\lambda' > \mu'$ then $c_\lambda = 0$. Therefore, if $T \in \Syt(\lambda)$ has $\Des(T) = \Des(T_\mu)$ and $c_\lambda \ne 0$, then $T = T_\mu$. Consequently, substituting $J = \Des(T_\mu)$ into Equation~\eqref{eqn:proof symmetric set with no consecutive indices is schur positive - specific set}, we find that $|\A(D = \Des(T_\mu))| = c_\mu \neq 0$.

    We may conclude that there exists an element $a \in \A$ with $D(a) = \Des(T_\mu)$. The set $D(a)$ is sparse, so $\{1,2\} \nsubseteq D(a)$. Lemma~\ref{lem:superstandard characterize coefficients} implies that $\{1,2\} \subseteq \Des(T_\mu)$ whenever $\mu_1' >2$, so we may deduce that $\mu_1' \le 2$. Therefore, $c_\lambda = 0$ for every partition $\lambda$ with more than $2$ parts.
    
    Thus, we can reformulate Equation~\eqref{eqn:proof symmetric set with no consecutive indices is schur positive - specific set} as
    \begin{equation}\label{eqn:proof symmetric set with no consecutive indices is schur positive:equation with only two rows}
        |\A(D = J)| = \sum_{k=0}^n c_k \left|\Syt(N-k,k)(\Des = J)\right|.
    \end{equation}
    
    By Lemma~\ref{lem:superstandard two rows characterize coefficients}, since only partitions into at most two parts are involved in Equation~\eqref{eqn:proof symmetric set with no consecutive indices is schur positive:equation with only two rows}, substituting $J = \Des(T_{\lambda})$ for $\lambda = (N-k,k)$ yields $|\A(D = \odds(k))| = c_k$, as required.
\end{proof}

\end{document}
